\documentclass[11pt,oneside]{book}

\usepackage[margin=1in]{geometry}
\usepackage{amsmath,amssymb}
\usepackage{graphicx}
\usepackage{hyperref}
\usepackage{bookmark}
\usepackage{microtype}
\usepackage{enumitem}
\usepackage{xcolor}

\hypersetup{
  colorlinks=true,
  linkcolor=blue,
  citecolor=blue,
  urlcolor=blue
}

% Simple placeholders helpers
\newcommand{\placeholder}[1]{%
  \par\noindent\textcolor{gray}{\textit{[Placeholder: #1]}}\par
}

\title{Decision Support Systems: A Modern Approach}
\author{Dr.\ Haitham El-Ghareeb\\
Associate Professor of Information Systems,\\
Faculty of Computers and Information Sciences,\\
Mansoura University, Egypt}
\date{\today}

\begin{document}

\frontmatter
\maketitle
\tableofcontents

\chapter*{Preface}
\addcontentsline{toc}{chapter}{Preface}
This book was prepared as a teaching text for a Decision Support Systems (DSS)
course with a practical orientation toward analytics, machine learning, and
organizational decision making. The aim is to bridge classical DSS foundations
with modern AI-enabled practice while keeping the material suitable for course
delivery, assignments, revision, and slide-based lecturing.

The sequence starts from core decision-making concepts, moves through DSS
architecture and model management, and then develops the analytics side of the
course through machine learning, business intelligence, expert systems,
unsupervised methods, and multi-criteria decision support. The final chapters
support revision through essay questions and a compact chapter of acronyms and
definitions.

\chapter*{How to Use This Book}
\addcontentsline{toc}{chapter}{How to Use This Book}
The chapters are designed to work at three levels:
\begin{itemize}
  \item \textbf{Conceptual study}: read the overview, main sections, and summary.
  \item \textbf{Assessment preparation}: use the key terms, MCQs, essay chapter,
  and midterm guidance.
  \item \textbf{Project development}: adapt the suggested DSS projects at the end
  of most chapters.
\end{itemize}

For a standard course sequence, Chapters 1--5 establish foundations, Chapters
7--12 deepen analytical and design capabilities, and Chapters 13--14 support
revision and consolidation.

\mainmatter

% ==========================================================
% Chapters (matching your course-topic list)
% ==========================================================

\include{chapters/ch01-introduction}

\include{chapters/ch02-decision-making-concepts}

\include{chapters/ch03-technologies-environment}

\include{chapters/ch04-model-management}

\include{chapters/ch05-intro-ml-in-dss}


\chapter{Midterm Review and Sample Questions}
\label{ch:midterm-review-and-sample-questions}
\section*{Midterm Overview}
\addcontentsline{toc}{section}{Midterm Overview}

The midterm is intended to evaluate the student’s understanding of the
foundational concepts developed in the first half of the course. Coverage
typically includes Chapters 1--5 together with the accompanying classroom
discussion and examples.

\subsection*{Recommended format}
\addcontentsline{toc}{subsection}{Recommended format}
\begin{itemize}
  \item \textbf{Part A:} short-answer and definition questions on DSS concepts,
  architecture, and decision-making terminology.
  \item \textbf{Part B:} analytical questions that compare models, evaluation
  approaches, and deployment choices.
  \item \textbf{Part C:} one short essay or case-based design question requiring
  a justified DSS design decision.
\end{itemize}

\subsection*{Preparation guidance}
\addcontentsline{toc}{subsection}{Preparation guidance}
\begin{itemize}
  \item Revise key distinctions such as MIS vs.\ BI vs.\ DSS and structured vs.\
  semi-structured vs.\ unstructured decisions.
  \item Be able to explain bounded rationality, common cognitive biases, and why
  governance matters in DSS.
  \item Review model management vocabulary including verification, validation,
  calibration, and data leakage.
  \item Practice describing how a prediction becomes an operational decision.
\end{itemize}

\section*{Sample Questions}
\addcontentsline{toc}{section}{Sample Questions}

\begin{enumerate}[label=\textbf{Q\arabic*.}]
  \item \textbf{Define a DSS and explain how it differs from a traditional MIS.}

  \textit{Brief answer guideline:} A strong answer should mention the focus on
  semi-structured and unstructured decisions, the role of models and interactive
  analysis, and the difference between reporting and decision support.

  \item \textbf{Why is bounded rationality important when designing a DSS
  interface?}

  \textit{Brief answer guideline:} A strong answer should connect bounded
  rationality to limited time, limited cognition, progressive disclosure,
  summaries, and bias mitigation.

  \item \textbf{Explain the difference between verification and validation in
  model management.}

  \textit{Brief answer guideline:} Verification asks whether the model was built
  correctly; validation asks whether the model is appropriate for the actual
  decision context.

  \item \textbf{A hospital wants to use an ML model to prioritize emergency
  cases. Which evaluation issues matter more than accuracy alone?}

  \textit{Brief answer guideline:} Expected points include recall for critical
  cases, calibration, threshold selection, operational capacity, fairness, and
  human review.
\end{enumerate}


\include{chapters/ch07-training-classification}

\include{chapters/ch08-supervised-learning}

\include{chapters/ch09-automated-expert-systems}

\include{chapters/ch10-business-intelligence}

\chapter{Unsupervised Learning and Recommender Systems for DSS}
\label{ch:unsupervised-learning-and-recommender-systems-for-dss}

\section{Chapter Overview and Learning Objectives}

Many decision-support problems do not begin with labeled outcomes. Instead,
they begin with raw operational, textual, or behavioral data where the system
must first discover structure before it can recommend an action. Unsupervised
learning methods help DSS identify segments, anomalies, latent topics, and
behavioral similarities, while recommender systems translate those patterns
into ranked options for users or managers.

\subsection*{Learning objectives}
\addcontentsline{toc}{subsection}{Learning objectives}
By the end of this chapter, the student should be able to:
\begin{itemize}
  \item distinguish unsupervised learning from supervised learning in DSS,
  \item explain how clustering, dimensionality reduction, and association
  analysis support decision making,
  \item describe the architecture of a recommender system in a DSS context,
  \item compare content-based, collaborative, and hybrid recommendation
  strategies,
  \item identify evaluation criteria for unsupervised and recommender systems,
  \item discuss risks such as popularity bias, cold start, and feedback loops.
\end{itemize}

\section{Unsupervised Learning in DSS}

\subsection{Why unlabeled structure matters}
In practice, organizations often have abundant data but few reliable labels.
Unsupervised learning helps DSS answer questions such as:
\begin{itemize}
  \item Which customers or students form meaningful segments?
  \item Which records look unusual enough to review?
  \item Which variables summarize the main structure of the data?
\end{itemize}

\subsection{Clustering for segmentation}
Clustering groups similar records without predefined classes. In DSS,
clustering is often used for:
\begin{itemize}
  \item customer segmentation for differentiated interventions,
  \item patient grouping for care pathways,
  \item student grouping for learning support,
  \item incident grouping for operational triage.
\end{itemize}
Common methods include $k$-means, hierarchical clustering, and density-based
methods such as DBSCAN.

\subsection{Dimensionality reduction}
High-dimensional data can be difficult to interpret and visualize.
Dimensionality reduction methods such as Principal Component Analysis (PCA) and
manifold-learning techniques help summarize the dominant structure of the data.
In DSS, this can support exploratory analysis, dashboard simplification, and
feature compression before downstream modeling.

\subsection{Association analysis}
Association rules discover co-occurrence patterns such as items frequently
purchased together, services used in sequence, or symptoms commonly observed
together. While association analysis does not by itself prescribe action, it
can drive bundling, cross-selling, fraud review, and knowledge discovery.

\section{From Patterns to Decisions}

\subsection{Anomaly and novelty detection}
Anomaly detection is useful when the objective is to identify cases that merit
review rather than produce a final automated decision. Examples include:
\begin{itemize}
  \item suspicious transactions,
  \item unusual patient deterioration signals,
  \item abnormal network events,
  \item unexpected operational bottlenecks.
\end{itemize}
The DSS should treat anomaly scores as prioritization signals and pair them
with human review workflows.

\subsection{Interpretation challenges}
Unsupervised models can produce patterns that are mathematically coherent but
operationally meaningless. For this reason, domain validation is essential:
\begin{itemize}
  \item inspect cluster profiles,
  \item test stability across time and samples,
  \item confirm that discovered patterns support actual decisions.
\end{itemize}

\section{Recommender Systems}

\subsection{What a recommender system does}
A recommender system ranks items, actions, or interventions for a user, group,
or case. In DSS, recommendation may target:
\begin{itemize}
  \item products or services,
  \item learning resources,
  \item treatment pathways,
  \item support actions for at-risk cases.
\end{itemize}

\subsection{Content-based recommendation}
Content-based recommenders rely on attributes of items and profiles of users or
cases. They are useful when rich metadata is available and when explaining the
recommendation in terms of item attributes is important.

\subsection{Collaborative filtering}
Collaborative filtering uses similarity in historical behavior across users and
items. It can reveal patterns not obvious from metadata alone, but it suffers
from cold-start problems when new users or new items have little interaction
history.

\subsection{Hybrid recommendation}
Many production DSS use hybrid recommenders that combine content features,
interaction history, business rules, and constraints. Hybrid systems are often
more robust and better aligned with policy requirements.

\section{Evaluation and Governance}

\subsection{Evaluating unsupervised methods}
Unsupervised methods are evaluated differently from supervised models. Typical
criteria include:
\begin{itemize}
  \item internal structure metrics such as silhouette score,
  \item stability across runs or time periods,
  \item interpretability and usefulness to decision makers,
  \item downstream impact on decisions and interventions.
\end{itemize}

\subsection{Evaluating recommenders}
Recommendation quality can be assessed using:
\begin{itemize}
  \item ranking metrics such as precision at $k$ and recall at $k$,
  \item online response metrics such as click-through or acceptance rate,
  \item diversity and novelty of recommendations,
  \item fairness and exposure across groups or items.
\end{itemize}

\subsection{Risks and limitations}
Important governance issues include:
\begin{itemize}
  \item \textbf{Cold start}: insufficient history for new users or items.
  \item \textbf{Popularity bias}: over-exposure of already popular items.
  \item \textbf{Feedback loops}: recommendations shape future data.
  \item \textbf{Filter bubbles}: excessive narrowing of user options.
  \item \textbf{Opaque segments}: clusters that are hard to justify in practice.
\end{itemize}

\section{Summary and Concluding Remarks}

This chapter showed how unlabeled data can still support decision making
through clustering, dimensionality reduction, association analysis, anomaly
detection, and recommendation. The main design principle is to connect
discovered structure to concrete actions: segmentation must lead to
interventions, anomaly scores must feed review queues, and recommendations must
respect organizational constraints and fairness requirements.

\section{Key Terms}
\begin{description}[style=nextline]
  \item[Unsupervised learning] Learning patterns from unlabeled data.
  \item[Clustering] Grouping similar observations without predefined labels.
  \item[Dimensionality reduction] Compressing high-dimensional data into a
  smaller representation while preserving important structure.
  \item[Association rule] A pattern describing items or events that frequently
  occur together.
  \item[Anomaly detection] Identifying unusual records or behaviors that may
  require review.
  \item[Recommender system] A system that ranks items or actions for a user or
  case.
  \item[Collaborative filtering] Recommendation using shared behavioral
  patterns across users and items.
  \item[Content-based recommendation] Recommendation using item and user
  attributes.
  \item[Cold start] The difficulty of recommending for new users or items with
  limited history.
  \item[Popularity bias] Over-recommending already popular items.
\end{description}

\section{Multiple-Choice Questions (MCQs)}

\subsection*{Questions}
\addcontentsline{toc}{subsection}{Questions}
\begin{enumerate}[label=\textbf{Q\arabic*.}]
  \item Clustering is most appropriate when:
  \begin{enumerate}[label=\alph*)]
    \item class labels are already known
    \item the DSS needs to discover latent segments in unlabeled data
    \item the objective is to solve a linear program
    \item only deterministic rules are allowed
  \end{enumerate}

  \item A common cold-start problem occurs when:
  \begin{enumerate}[label=\alph*)]
    \item there is too much labeled data
    \item a recommender has little history for new users or items
    \item clusters overlap perfectly
    \item dimensionality is low
  \end{enumerate}

  \item Which issue is most associated with recommender-system feedback loops?
  \begin{enumerate}[label=\alph*)]
    \item recommendations change future behavior and future data
    \item the system cannot rank items
    \item rule-based reasoning becomes impossible
    \item all items receive equal exposure
  \end{enumerate}

  \item PCA is primarily used for:
  \begin{enumerate}[label=\alph*)]
    \item human approval workflows
    \item dimensionality reduction and structure summarization
    \item legal compliance audits
    \item backward chaining
  \end{enumerate}
\end{enumerate}

\subsection*{Answer Key}
\addcontentsline{toc}{subsection}{Answer Key}
\begin{enumerate}[label=\textbf{Q\arabic*.}]
  \item b
  \item b
  \item a
  \item b
\end{enumerate}


\chapter{Multi-Criteria Decision Analysis for DSS}
\label{ch:multi-criteria-decision-analysis-for-dss}

\section{Chapter Overview and Learning Objectives}

Many important DSS problems cannot be reduced to a single objective. Managers
often balance cost, quality, risk, fairness, time, and strategic fit at the
same time. Multi-Criteria Decision Analysis (MCDA) provides structured methods
for evaluating alternatives when several criteria matter simultaneously.

\subsection*{Learning objectives}
\addcontentsline{toc}{subsection}{Learning objectives}
By the end of this chapter, the student should be able to:
\begin{itemize}
  \item explain why MCDA is important in semi-structured DSS problems,
  \item construct a criteria matrix and normalize decision criteria,
  \item compare weighted scoring, AHP, and outranking-style reasoning at a high
  level,
  \item perform sensitivity analysis on criteria weights,
  \item identify governance risks in criteria design and weighting.
\end{itemize}

\section{Why Multi-Criteria Methods Matter}

\subsection{Single-objective models are often insufficient}
Real DSS problems often involve competing objectives:
\begin{itemize}
  \item selecting a supplier with trade-offs among cost, quality, and delivery,
  \item prioritizing projects under risk, benefit, and strategic alignment,
  \item ranking interventions where effectiveness, equity, and feasibility all
  matter.
\end{itemize}

\subsection{Core MCDA workflow}
A practical MCDA workflow in DSS includes:
\begin{enumerate}[label=\arabic*.]
  \item define alternatives,
  \item define criteria and measurement scales,
  \item determine directions of preference,
  \item assign or elicit weights,
  \item aggregate or compare alternatives,
  \item test sensitivity and robustness,
  \item document rationale and final decision.
\end{enumerate}

\section{Weighted Scoring Models}

\subsection{Decision matrix}
The starting point is a decision matrix where rows represent alternatives and
columns represent criteria. Because criteria can have different units,
normalization is often required before aggregation.

\subsection{Weighted sum model}
One of the simplest approaches is the weighted sum:
$$
S(a) = \sum_{j=1}^{m} w_j x_{aj},
$$
where $w_j$ is the weight of criterion $j$ and $x_{aj}$ is the normalized score
of alternative $a$ on criterion $j$.

\subsection{Strengths and limitations}
Weighted scoring is transparent and easy to implement in DSS dashboards, but it
assumes compensability: poor performance on one criterion can be offset by high
performance on another. In some settings, this is undesirable.

\section{Analytic Hierarchy Process (AHP)}

\subsection{Hierarchical structuring}
AHP decomposes the problem into goal, criteria, subcriteria, and alternatives.
Decision makers compare criteria pairwise, which can be more intuitive than
assigning direct numeric weights.

\subsection{Consistency}
An advantage of AHP is that it checks the consistency of pairwise judgments. If
judgments are highly inconsistent, the DSS can prompt the user to revise them
before ranking alternatives.

\section{Sensitivity and Robustness}

\subsection{Why sensitivity analysis is essential}
Small changes in weights can reverse rankings. A DSS should therefore expose:
\begin{itemize}
  \item how rankings change when weights change,
  \item which criteria dominate the result,
  \item where alternatives are nearly tied.
\end{itemize}

\subsection{Scenario-based MCDA}
One useful pattern is to define stakeholder-specific or scenario-specific
weighting schemes. For example, a finance perspective may weight cost more
heavily, while a service perspective may weight quality and responsiveness more
heavily.

\section{Governance Considerations}

\subsection{Criteria design is a policy decision}
MCDA can appear objective while embedding subjective choices. Governance
therefore requires transparent documentation of:
\begin{itemize}
  \item why criteria were selected,
  \item how weights were chosen,
  \item who approved the weighting scheme,
  \item how conflicts are resolved.
\end{itemize}

\subsection{Fairness and accountability}
If criteria disadvantage specific groups, the DSS may institutionalize bias.
Criteria and weights should therefore be reviewed for fairness, legality, and
organizational legitimacy.

\section{Summary and Concluding Remarks}

MCDA provides a disciplined way to evaluate alternatives when several criteria
must be balanced. In DSS practice, its value is not only computational but also
communicative: it makes trade-offs explicit, supports stakeholder discussion,
and documents why an alternative was ranked highly. The most important design
principle is transparency in criteria, weights, and sensitivity analysis.

\section{Key Terms}
\begin{description}[style=nextline]
  \item[MCDA] Multi-Criteria Decision Analysis; a family of methods for
  evaluating alternatives across several criteria.
  \item[Decision matrix] Table of alternatives and their scores on criteria.
  \item[Normalization] Transforming criteria values to comparable scales.
  \item[Weight] Relative importance assigned to a criterion.
  \item[Compensability] The extent to which poor performance on one criterion
  can be offset by strong performance on another.
  \item[AHP] Analytic Hierarchy Process; a structured MCDA method using pairwise
  comparisons.
  \item[Sensitivity analysis] Studying how changes in assumptions affect the
  ranking or recommendation.
\end{description}

\section{Multiple-Choice Questions (MCQs)}

\subsection*{Questions}
\addcontentsline{toc}{subsection}{Questions}
\begin{enumerate}[label=\textbf{Q\arabic*.}]
  \item MCDA is especially useful when:
  \begin{enumerate}[label=\alph*)]
    \item there is only one criterion and no constraints
    \item decision makers must balance multiple competing criteria
    \item labels are unavailable for clustering
    \item a rule base replaces all judgment
  \end{enumerate}

  \item A weighted scoring model requires:
  \begin{enumerate}[label=\alph*)]
    \item criteria, scores, and weights
    \item only a confusion matrix
    \item only symbolic rules
    \item no normalization at all
  \end{enumerate}

  \item Sensitivity analysis helps determine:
  \begin{enumerate}[label=\alph*)]
    \item whether rankings are robust to weight changes
    \item whether a database is normalized
    \item whether an image classifier is deep enough
    \item whether all alternatives have equal cost
  \end{enumerate}

  \item AHP is known for:
  \begin{enumerate}[label=\alph*)]
    \item pairwise comparison of criteria and consistency checking
    \item only using random selection
    \item ignoring stakeholder judgment
    \item replacing all decision matrices
  \end{enumerate}
\end{enumerate}

\subsection*{Answer Key}
\addcontentsline{toc}{subsection}{Answer Key}
\begin{enumerate}[label=\textbf{Q\arabic*.}]
  \item b
  \item a
  \item a
  \item a
\end{enumerate}


\chapter{Essay Questions with Model Answers}
\label{ch:essay-questions-with-model-answers}

\section{How to Use This Chapter}

This chapter is intended for revision, exam preparation, and assignment design.
Each essay question is followed by a concise model answer outline. The answers
are not the only acceptable responses; rather, they illustrate the depth,
structure, and terminology expected in a strong DSS course answer.

\section{Essay Questions and Model Answers}

\subsection*{Essay 1: Explain the difference between MIS, BI, and DSS.}
\addcontentsline{toc}{subsection}{Essay 1}
\textbf{Model answer.}
Management Information Systems (MIS) mainly support routine and structured
reporting. Business Intelligence (BI) extends this by integrating data,
defining KPIs, and presenting descriptive and diagnostic analytics through
dashboards and reports. Decision Support Systems (DSS) go further by supporting
semi-structured and unstructured decisions through models, scenarios,
recommendations, and interactive analysis. A modern DSS may use BI as its data
foundation while adding prediction, optimization, rules, and human-in-the-loop
workflows.

\subsection*{Essay 2: Why is human-in-the-loop design important in AI-enabled DSS?}
\addcontentsline{toc}{subsection}{Essay 2}
\textbf{Model answer.}
Human-in-the-loop design is important because many DSS decisions are high
impact, uncertain, and context dependent. AI models can produce useful
predictions, but they may also be wrong, biased, poorly calibrated, or unable
to represent policy constraints. Human review provides accountability, domain
judgment, and exception handling. In well-designed DSS, the system computes and
prioritizes, while the human validates, overrides when needed, and remains
responsible for final action.

\subsection*{Essay 3: Discuss the role of model management in DSS.}
\addcontentsline{toc}{subsection}{Essay 3}
\textbf{Model answer.}
Model management is the discipline of organizing, validating, deploying,
monitoring, and governing the models used in DSS. It covers optimization
models, simulation models, statistical forecasts, and ML/DL models. Good model
management ensures that models are versioned, documented, validated against the
intended decision context, monitored for drift, and linked to appropriate
approval workflows. Without model management, organizations risk using outdated
or unreliable models in important decisions.

\subsection*{Essay 4: Compare supervised and unsupervised learning in DSS.}
\addcontentsline{toc}{subsection}{Essay 4}
\textbf{Model answer.}
Supervised learning uses labeled data to predict known outcomes such as class
labels or numeric targets. In DSS, it supports risk scoring, forecasting, and
classification. Unsupervised learning works without labels and instead
discovers structure such as clusters, latent dimensions, associations, or
anomalies. In DSS, it supports segmentation, exploratory analysis, anomaly
review, and recommendation. Supervised learning is typically evaluated through
predictive performance metrics, while unsupervised methods require stability,
interpretability, and downstream decision usefulness.

\subsection*{Essay 5: How do expert systems differ from machine-learning systems?}
\addcontentsline{toc}{subsection}{Essay 5}
\textbf{Model answer.}
Expert systems represent knowledge explicitly through facts, rules, and an
inference engine. They are transparent and useful for encoding policy,
compliance, and expert judgment. Machine-learning systems infer patterns from
data and are often better at handling complex relationships and large-scale
prediction. Expert systems are easier to audit but may become brittle when the
environment changes. ML systems are adaptive but may be less interpretable.
Modern DSS increasingly combine both approaches in hybrid systems where ML
produces scores and rules enforce guardrails.

\subsection*{Essay 6: Explain the relationship between BI and DSS.}
\addcontentsline{toc}{subsection}{Essay 6}
\textbf{Model answer.}
BI and DSS are closely related but not identical. BI focuses on collecting,
integrating, governing, and visualizing data so that organizational performance
is visible and measurable. DSS uses that information, often together with
models and expert knowledge, to evaluate alternatives and recommend actions.
BI answers questions such as ``what happened'' and ``why did it happen,'' while
DSS increasingly addresses ``what will happen'' and ``what should we do.'' In
practice, a robust DSS usually depends on strong BI foundations.

\subsection*{Essay 7: Why must fairness, privacy, and accountability be addressed in DSS?}
\addcontentsline{toc}{subsection}{Essay 7}
\textbf{Model answer.}
DSS can materially affect people, resources, and institutional outcomes.
Fairness matters because biased data, criteria, or models can systematically
harm particular groups. Privacy matters because decision support often uses
sensitive personal or organizational data. Accountability matters because
decisions must be explainable, reviewable, and contestable, especially in
high-stakes settings. Ethical DSS design therefore requires governance,
documentation, audit trails, and clearly assigned human responsibility.

\subsection*{Essay 8: What is the value of sensitivity analysis in DSS?}
\addcontentsline{toc}{subsection}{Essay 8}
\textbf{Model answer.}
Sensitivity analysis studies how a recommendation changes when assumptions,
inputs, thresholds, or weights change. It is valuable because many DSS outputs
depend on uncertain estimates, subjective preferences, or unstable data.
Sensitivity analysis helps identify which factors drive the recommendation,
whether the decision is robust, and where additional information would be most
useful. It improves both technical confidence and stakeholder trust.

\section{Instructor Notes}

These model answers can be expanded into long-form assignments, oral
examination prompts, or take-home essay questions. In assessment use, students
should be expected to provide examples, connect concepts across chapters, and
argue design trade-offs rather than merely reproduce definitions.


\chapter{Acronyms and Definitions}
\label{ch:acronyms-and-definitions}

\section{Acronyms}

\begin{description}[style=nextline]
  \item[AI] Artificial Intelligence
  \item[AHP] Analytic Hierarchy Process
  \item[API] Application Programming Interface
  \item[BI] Business Intelligence
  \item[CDC] Change Data Capture
  \item[CDSS] Clinical Decision Support System
  \item[CNN] Convolutional Neural Network
  \item[CRM] Customer Relationship Management
  \item[DBSCAN] Density-Based Spatial Clustering of Applications with Noise
  \item[DL] Deep Learning
  \item[DNN] Deep Neural Network
  \item[DSS] Decision Support System
  \item[ELT] Extract, Load, Transform
  \item[ERP] Enterprise Resource Planning
  \item[ETL] Extract, Transform, Load
  \item[GDSS] Group Decision Support System
  \item[GPU] Graphics Processing Unit
  \item[IoT] Internet of Things
  \item[KPI] Key Performance Indicator
  \item[LSTM] Long Short-Term Memory
  \item[MCDA] Multi-Criteria Decision Analysis
  \item[MIS] Management Information System
  \item[ML] Machine Learning
  \item[NLP] Natural Language Processing
  \item[OLAP] Online Analytical Processing
  \item[OLTP] Online Transaction Processing
  \item[PCA] Principal Component Analysis
  \item[RNN] Recurrent Neural Network
  \item[SVM] Support Vector Machine
  \item[XAI] Explainable Artificial Intelligence
\end{description}

\section{Core Definitions}

\begin{description}[style=nextline]
  \item[Decision Support System (DSS)] A computer-based system that supports
  semi-structured and unstructured decision making through data, models,
  knowledge, and interactive interfaces.
  \item[Structured decision] A routine decision with a well-defined procedure.
  \item[Semi-structured decision] A decision that combines formal analysis with
  human judgment.
  \item[Unstructured decision] A novel or complex decision with no complete
  predefined procedure.
  \item[Model management] The organization, validation, deployment, monitoring,
  and governance of analytical models used by a DSS.
  \item[Business Intelligence] Technologies and processes used to integrate,
  analyze, and visualize data for decision making.
  \item[Expert system] A knowledge-driven system using explicit rules and an
  inference engine to provide recommendations.
  \item[Human-in-the-loop] A design pattern in which humans review, approve, or
  override algorithmic recommendations.
  \item[Calibration] The degree to which predicted probabilities match observed
  outcome frequencies.
  \item[Data leakage] Use of information during model development that would not
  be available at decision time.
  \item[Cold start] A recommendation problem caused by limited history for new
  users or items.
  \item[Sensitivity analysis] Examination of how outputs change when inputs,
  assumptions, or weights are varied.
  \item[Audit trail] A documented record of data, model/rule versions, and
  actions used for accountability.
  \item[Fairness] The requirement that DSS decisions and recommendations avoid
  unjustified systematic disadvantage across groups.
  \item[Explainability] The degree to which a human can understand why a system
  produced a particular output or recommendation.
\end{description}

\section{How to Use This Chapter}

This chapter is intended as a compact revision aid. Students can use it before
exams, during slide review, and while preparing projects to ensure that core
terminology remains precise and consistent across the course.


\backmatter
\chapter*{Appendix A: Python Setup and Tooling}
\addcontentsline{toc}{chapter}{Appendix A: Python Setup and Tooling}
\placeholder{Environment setup, packages, and reproducibility.}

\chapter*{Appendix B: Mathematical Background}
\addcontentsline{toc}{chapter}{Appendix B: Mathematical Background}
\placeholder{Probability, linear algebra, optimization basics.}

\chapter*{Bibliography}
\addcontentsline{toc}{chapter}{Bibliography}
\placeholder{Add Bib\TeX{} file later.}

\chapter*{Index}
\addcontentsline{toc}{chapter}{Index}
\placeholder{Index placeholder.}

\end{document}
